% Generated draft: requires mathematical and source review.
\documentclass[11pt]{article}
\usepackage[margin=25mm]{geometry}
\usepackage{fontspec}
\setmainfont{Latin Modern Roman}
\usepackage{amsmath,amssymb,mathtools,unicode-math}
\setmathfont{Latin Modern Math}
\usepackage{enumitem,needspace,xurl,hyperref}
\setlength{\parindent}{0pt}
\setlength{\parskip}{5pt}
\setlength{\emergencystretch}{3em}
\newcommand{\N}{\mathbb{N}}
\newcommand{\Z}{\mathbb{Z}}
\newcommand{\Q}{\mathbb{Q}}
\newcommand{\R}{\mathbb{R}}
\newcommand{\C}{\mathbb{C}}
\newcommand{\F}{\mathbb{F}}
\newcommand{\A}{\mathbb{A}}
\newcommand{\PP}{\mathbb{P}}
\newcommand{\E}{\mathbb{E}}
\newcommand{\eps}{\varepsilon}
\newcommand{\dd}{\,\mathrm{d}}
\begin{document}
\section*{UnsolvedMath ID: 701747}
\textbf{Source problem code:} catalogue entry
\subsection{Definitions and mathematical statement}
Let \(K\) be a field, let \(d\geq 1\) be an integer, and let \(G\leq \operatorname{GL}_d(K)\). This uses the usual finite-dimensional meaning of a linear group; the supplied source specifies no restriction on the field or its characteristic. Use the commutator convention \([x,g]=x^{-1}g^{-1}xg\), and define the left-normed iterated commutators by
\[
c_1(x,g)=[x,g],\qquad c_{n+1}(x,g)=[c_n(x,g),g]\quad(n\geq 1).
\]
An element \(g\in G\) is \emph{almost Engel} if
\[
\exists E_g\subseteq G\text{ finite}\quad
\forall x\in G\quad
\exists N=N(x,g)\in\Z_{\geq 1}\quad
\forall n\geq N,\qquad c_n(x,g)\in E_g.
\]
Thus the finite set \(E_g\) may depend on \(g\), but must work for every \(x\); the threshold may depend on both \(x\) and \(g\). An element is \emph{Engel} if the same condition holds with \(E_g=\{1_G\}\). Put
\[
\operatorname{AE}(G)=\{g\in G:g\text{ is almost Engel}\}.
\]
\textbf{Question (P. Shumyatsky, Kourovka Notebook Problem 20.89).} For every field \(K\), every integer \(d\geq 1\), and every subgroup \(G\leq\operatorname{GL}_d(K)\), is \(\operatorname{AE}(G)\) a subgroup of \(G\)?

The supplied record gives two background results: Gruenberg's result that the Engel elements of a linear group form a subgroup, and Shumyatsky's result that a linear group all of whose elements are almost Engel is finite-by-hypercentral. Here finite-by-hypercentral means that there is a finite normal subgroup with hypercentral quotient, that is, a quotient equal to its transfinite upper center. These are supplied contextual assertions, not additional hypotheses or independently verified results.
\subsection{Short English statement}
Does the set of almost Engel elements form a subgroup in every linear group? For each such element \(g\), one finite set must eventually contain all sufficiently long iterated commutators with \(g\), for every starting element \(x\), with the threshold allowed to depend on \(x\) and \(g\). The supplied registry labels the question open; this draft does not independently verify that status.
\subsection{Sources}
\begin{itemize}
\item Local UnsolvedMath registry, permanent ID \(701747\), exact source code \texttt{catalogue entry}, title \emph{Kourovka Notebook Problem 20.89}. Supplied upstream revision: \texttt{d371ee917ac9647fbb06cd6fc529deec6f3dda82}. The record was updated on 2026-10-16 and supplies the status \emph{open}.
\item \emph{The Kourovka Notebook (Unsolved Problems in Group Theory)}, 20th issue (2022), Problem 20.89, proposed by P. Shumyatsky. The supplied background also identifies the 21st issue edition (2026), arXiv:1401.0300v48. The collection description names E. I. Khukhro and V. D. Mazurov as editors.
\item Gruenberg, \emph{J. Algebra} \textbf{3} (1966), 291--303. Supplied citation for the assertion that the Engel elements of a linear group form a subgroup; no article title or link is supplied.
\item P. Shumyatsky, \emph{Monatsh. Math.} \textbf{186} (2018), 711--719. Supplied citation for the assertion that a linear group in which every element is almost Engel is finite-by-hypercentral; no article title or link is supplied.
\end{itemize}
\end{document}
